\documentclass[11pt]{article}
\usepackage[a4paper,margin=26mm]{geometry}
\usepackage{amsmath,amssymb,amsthm,booktabs,hyperref}
\newtheorem{theorem}{Theorem}\newtheorem{lemma}[theorem]{Lemma}
\newcommand{\Aut}{\operatorname{Aut}}
\title{Primitive geometries from subgroup pairs}
\author{}\date{}
\begin{document}\maketitle
We treat the mathematical results in the twelve-page journal edition of
Dimitri Leemans and Bernardo G. Rodrigues, \emph{On primitive geometries
of rank two}, Art Discrete Appl. Math. 2 (2019), P1.08. The original
Rudvalis count is false. The previously accepted orbit bound already
proved this; the present treatment supplies the exact corrected count,
recovers the other eleven counts, and accounts for every finite table in
the article. It also records a general algorithm correction, two exchanged
subgroup labels and four exchanged pairs of type diameters.

This is a treatment of that fixed article. It does not certify the contents
of the electronic representative files linked from its page 2, nor claim
identity with every edition of those external datasets. The distinction
between complete mathematical parametrization and file verification is
made precise below. Historical running times remain attributed observations.

\section{The common classification and counting argument}
We first recall the elementary construction attributed by the source to
Tits. Given subgroups $G_i\leq G$, take the right cosets $G_i g$ as objects
of type $i$, with incidence defined by nonempty intersection. The base
cosets $G_i$ form a chamber since all contain identity. Right multiplication
preserves types and incidence. It is transitive on objects of a given type,
and on incident pairs of any two fixed types: a common element $g$ sends
the base pair to the given pair. These statements concern the incidence
structure and flags of size at most two; no general higher-rank
flag-transitivity assertion is being assumed.
Cosets are tagged by their types, so two equal subsets with different
types remain different objects.

For distinct maximal subgroups $H,K$ in a finite group, the rank-two
incidence graph is connected because $\langle H,K\rangle=G$. Its two
degrees are $[H:H\cap K]$ and $[K:H\cap K]$, both at least two. Maximality
makes the two coset actions primitive. Thus the source's geometries are
precisely unordered pairs of distinct maximal subgroups, with equivalence
defined by $\Aut(G)$. The distinguished group action is retained; this is
not a claim about every abstract isomorphism of the underlying graphs.

\begin{theorem}\label{param}
Put $A=\Aut(G)$ and choose one representative $M_i$ of each $A$-orbit of
maximal subgroups of $G$. Let $N_i=N_A(M_i)$. A complete parameter set for
the source's primitive rank-two geometries is
\[
\coprod_{i<j}N_i\backslash A/N_j\quad\sqcup\quad
\coprod_i\Bigl((N_i\backslash A/N_i)\setminus\{N_i\}\Bigr)
                       /(N_i a N_i\sim N_i a^{-1}N_i).
\tag{1}
\]
The parameter represented by $a$ gives the coset geometry of
$\{M_i^a,M_j\}$. Here the automorphisms act on the right, so
$(M^a)^b=M^{ab}$; for an inner automorphism this is the familiar
notation $M^a=a^{-1}Ma$. In general $N_i$ means the setwise stabilizer
of $M_i$ in $A$.
\end{theorem}
\begin{proof}
Any pair of maximal subgroups can be written
$\{M_i^x,M_j^y\}$. Conjugating by $y^{-1}$ gives
$\{M_i^{xy^{-1}},M_j\}$. When $i\ne j$, the two distinct $A$-classes fix
the roles of the two subgroups. Two resulting pairs are equivalent if and
only if their parameters differ by left multiplication by $N_i$ and right
multiplication by $N_j$. This is exactly the double-coset relation.
When $i=j$, the double coset $N_i$ gives equal subgroups and must be
removed. Interchanging the two subgroups and then conjugating by $a^{-1}$
replaces $a$ by $a^{-1}$. This is the only additional identification.
The preceding coset argument supplies a firm, connected primitive geometry
for every remaining parameter.
\end{proof}

This description supplies every class and the exact equivalence relation,
not merely a lower bound. It is finite and constructive once the finite
group and its maximal representatives are supplied. One may list its
double cosets, or use the corrected coset-orbit algorithm proved in the
accompanying algorithm note. That note retains the original second
subgroup and lifts only the point stabilizer. For a general group, lifting
both images replaces $K$ by $K\operatorname{core}_G(H)$; a four-element
counterexample proves that the original unrestricted pseudocode is false.
For the twelve simple groups in the source, the proper coset actions are
faithful, so this correction does not itself change their totals.

To count (1), let $\Omega$ be the set of maximal subgroups and let
$\psi(a)=|\Omega^a|$. A fixed two-element subset either consists of two
fixed elements or is a two-cycle. Its number is therefore
\[
\frac{\psi(a)^2-2\psi(a)+\psi(a^2)}2.
\]
Burnside's lemma gives
\[
\#\binom{\Omega}{2}/A=
 \frac12\sum_{[a]_A}
 \frac{\psi(a)^2-2\psi(a)+\psi(a^2)}{|C_A(a)|}.             \tag{2}
\]
Here $\Omega$ is the disjoint union of the transitive actions $A/N_i$.
The accompanying common proof establishes the normalizer selection and
shows how complete character or mark data give every term of (2).
Its product-of-marks theorem also recovers each double-coset orbit-length
multiset. These are classical counting methods; their classical origin is
not presented as a new discovery.

\section{The twelve group-specific inputs}
For $M_{11},M_{23},M_{24},J_1,Ru$, the accepted character application uses
complete maximal-subgroup lists, their stored character fusions, centralizers
and square maps. The numbers of maximal classes are respectively
$5,7,9,7,15$. Since these groups have trivial outer automorphism groups,
$A=G$ and maximal subgroups are self-normalizing. Their primitive
permutation characters therefore give exactly the summands of $\psi$.
The named-group and complete-library contracts remain explicit premises.

For $M_{12},J_2,J_3,HS,McL,He$, the acquired complete subgroup tables of
$A=G.2$ identify the socle $G$, its maximal subgroups up to $A$, and their
actual normalizer classes in $A$. The resulting numbers of maximal
$A$-classes are $8,9,8,10,10,10$. Their normalizer mark rows, evaluated at
the stored cyclic-class fusion, give $\psi$ on every element class of $A$,
including the outer classes. Thus neither (1) nor (2) omits outer
identifications. The exact maximal and normalizer class lists are retained
in the group-specific input map, with their table editions.

For $M_{22}$, the previously accepted L073 census supplies the complete
$\Aut(M_{22})$ action and its $81$ rank-two classes. Its representative
construction and completeness proof are imported, not rerun. The separate
simple-group mark calculation gives $104$ inner classes, which must not be
substituted for the required full-automorphism total.

These inputs specify every family in (1), while (2) computes its exact
cardinality. They do not assert that permutation characters alone distinguish
all embeddings: the two degree-$253$ maximal subgroups of $M_{23}$ are a
concrete exception, resolved by the additional normal-subgroup certificate.
The explicit $M_{11}$ calculation further constructs all $37$ representatives,
their coset graphs and their design data.

\begin{center}\small
\begin{tabular}{lrrl}
\toprule
Group & printed total & recovered total & input route\\\midrule
$M_{11}$&37&37&characters and literal pairs\\
$M_{12}$&166&166&marks of $M_{12}.2$\\
$M_{22}$&81&81&accepted L073 census\\
$M_{23}$&170&170&characters\\
$M_{24}$&5074&5074&characters\\
$J_1$&669&669&characters\\
$J_2$&334&334&marks of $J_2.2$\\
$J_3$&546&546&marks of $J_3.2$\\
$HS$&339&339&marks of $HS.2$\\
$McL$&443&443&marks of $McL.2$\\
$He$&4074&4074&marks of $He.2$\\
$Ru$&1511&556287&characters\\\bottomrule
\end{tabular}
\end{center}
The earlier Rudvalis lower bound $50669$ already disproved $1511$.
The exact value $556287$ completes that correction; it is not a second
principal refutation. Summing the corrected table gives $568220$ classes
across the twelve groups, under the specified group equivalences.

The groups' standard orders and the degrees of the source's chosen
permutation representations are retained as named-group data. The displayed
runtimes concern the authors' implementation and workstation. Our exact
counting proof does not reproduce those benchmarks or certify a speedup.

\section{The finite tables and the correction ledger}
The product-of-marks calculation recovers all $308$ numerical cells in
Tables 2 and 5--8. In Table 7 the two degree-$253$ M23 row labels are
interchanged. The acquired class of order $40320$ with orbits $7,16$ on
the $23$ points contains a unique conjugate of an elementary abelian
subgroup of order $16$, and is $2^4:A_7$. The other is $L_3(4):2$.
The detailed proof counts contained conjugates by
$m_{M,L}|M|/|N_G(L)|$ and uses the complete cyclic fusion to prove exponent
two. No numeric table cell changes.

The preliminary upper-triangle counts $38$ for M11 and $268$ for M12 are
also recovered. Identifying diagonal double cosets under inversion gives the
inner counts $37$ and $255$; the full M12 automorphism action then gives
$166$. The three quantities have different meanings and are not compared
as if they were the same enumeration.

For Tables 3 and 4, the concrete M11 certificate covers $309$ maximal
subgroups and all $47586$ unordered pairs. It obtains $37$ orbits, checks
all coset graphs and transports every named intersection to an accepted
structural proof. All $70$ design parameters agree. With the printed
subgroup columns retained, four type-diameter triples change:
\[
\begin{array}{c|cc}
\text{source rows}&\text{printed}&\text{corrected}\\\hline
12,18 &[3,2,4]&[4,2,3]\\
23,35 &[3,3,4]&[4,3,3].
\end{array}
\]
The definition of $d_p$ and $d_l$ is in \emph{Section 2, physical page 3},
before Section 3 begins. The frozen M11 application mistakenly called this
Section 3. The present synthesis corrects that citation only; its source
page, four numerical corrections and proofs are unchanged. The publication
edition of the M11 text records the same explicit citation delta.

The M11 note gives a direct witness for row 35: the graph is the subdivision
of $K_{11}$, with eccentricity four at an edge vertex and three at an
original vertex. The two subgroup indices in the printed order are $55$
and $11$. This explains the correction independently of the graph program.
Rows 10/11 and 16/17 have indistinguishable printed profiles; both possible
matches are retained, while their actual pair orbits remain distinct.

The source omits design rows 31 and 32. They are complete incidence with
repeated full blocks. They meet its literal subset-counting definition,
but not a convention requiring proper or distinct blocks. We retain that
qualification explicitly. It does not change any of the seventy printed
design parameters. The undefined $K$ in the pseudocode is a separate
variable typo; it is distinct from the substantive core-kernel issue.

\section{Precisely what is, and is not, completed}
Theorem~\ref{param}, the corrected finite enumerator, the acquired complete
group-specific inputs and the exact orbit counts provide a complete
mathematical parametrization for all twelve groups. The individual finite
invariants printed in the paper are recovered with the corrections just
listed. The source's comparison to earlier M11, M22 and M23 totals is
accounted for by these same exact counts and the L073 import.
We do not expand (1) into a file listing all $568220$ representatives.
The classification supplied here is the explicit parameter set, its
equivalence relation and its exact size, with all the finite invariants
actually printed in the article checked separately.

The external data addresses on source page 2 are
\texttt{data.adam-journal.eu/apg/} and the authors' former
\texttt{homepages.ulb.ac.be/\string~dleemans/PRI/} directory. Their downloadable
representative files have not been read class by class, matched against
our parameters, or certified complete by this treatment. In particular,
our exact Ru count does not validate any external file carrying the old
$1511$ total. Those files are separate data artifacts, not extra input
premises silently used to prove the article's finite tables.

The Tits construction and the earlier Dehon algorithm are attributed
background results. The source's McLaughlin-code examples refer to another
paper; they are not newly proved code theorems here. Its suggestion to
extend the new algorithm to higher ranks is explicitly future work, and
its explanation of a possible earlier M22 error is speculation. Neither
is converted into a theorem or a solved-conjecture claim.

Our contribution is the common classification and counting treatment,
its applications with explicit input contracts, the finite certificates,
and the stated corrections. Burnside's lemma, tables of marks, transporter
counting and coset geometry are borrowed mathematics. The structural
classifications and library completeness used in the applications remain
visible premises. This source treatment claims no new ordinary recovery,
second principal refutation, Lean formalization or compression factor.
Acceptance of the complete corrected disposition is a separate review gate.
\end{document}
